\chapter{El Spillable Edge Store}
\label{ch:spill}

Ingestar un dump EVTX de 28~GB en una laptop con 16~GB de RAM
requiere volcar (\emph{spill}) las estructuras en memoria a disco
antes de que el proceso sea liquidado por OOM.
\code{SpillableEdgeStore} lo implementa para la lista de aristas:
las aristas se acumulan en un buffer en memoria hasta cruzar un
umbral; pasado el umbral, el buffer se ordena y se escribe como un
archivo ``run'' ordenado en disco, y se arranca un buffer nuevo.
En query time, el store presenta una única lista de adyacencia
lógica haciendo merge de todos los runs en \code{finalize}.

\section{Modos}

\begin{definition}[Modo del store]
El store está \emph{en memoria} (sin runs en disco) o
\emph{spilled} (uno o más runs en disco). Un tercer estado
``finalizado'' se alcanza después de que \code{finalize} unificó
todos los runs en un único archivo mmap; una vez finalizado, el
store es de sólo lectura.
\end{definition}

Campos:

\begin{lstlisting}[style=rust]
struct SpillableEdgeStore {
    buffer:       Vec<CompactRelation>,
    buffer_index: HashMap<StrId, Vec<u32>>, // posiciones en buffer
    runs:         Vec<SpillRunHandle>,      // runs ordenados en disco
    merged_mmap:  Option<MmapFile>,         // estado finalizado
    finalized:    bool,
    spill_limit:  usize,                    // aristas antes de flush
}
\end{lstlisting}

\code{buffer\_index[s]} guarda las \emph{posiciones} (índices u32)
de las aristas de la source \code{s} dentro del buffer actual, no
las aristas mismas. Esto hace barato a \code{get\_edges(s)} en
modo en memoria: el llamador obtiene un slice cortado por esas
posiciones, sin una copia HashMap-of-Vecs.

\section{Append}

\begin{algorithm}
\caption{\textsc{Spill-Append}}\label{alg:spill-append}
\KwIn{store $S$, source $s$, compact $c$}
\If{$S.\text{buffer.len()} \geq S.\text{spill\_limit}$}{
    \textsc{Flush-Buffer-To-Run}($S$)\Comment*[r]{ver abajo}
}
$p \gets S.\text{buffer.len()}$\;
\code{S.buffer.push($c$)}\;
\code{S.buffer\_index[$s$].push($p$)}\;
\end{algorithm}

\begin{algorithm}
\caption{\textsc{Flush-Buffer-To-Run}}\label{alg:spill-flush}
\KwIn{store $S$}
ordenar $S.\text{buffer}$ lexicográficamente por $(\text{source\_sid},
\text{timestamp}, \text{dest\_sid})$\;
escribir el buffer ordenado en un archivo temp $f$ como bytes
crudos (con header que lleva count y versión)\;
$S.\text{runs}.\text{push}(\text{handle}(f))$\;
vaciar $S.\text{buffer}$ y $S.\text{buffer\_index}$\;
\end{algorithm}

\begin{complexity}
\textsc{Spill-Append} es $\Ord{1}$ amortizado. Cada
\code{spill\_limit} appends dispara un \textsc{Flush} de costo
$\Ord{L \log L}$ donde $L$ = \code{spill\_limit}, distribuido a lo
largo del stream de ingesta. El costo total en modo spill para
ingerir $N$ aristas es
$\Ord{N \log L + N/L \cdot L} = \Ord{N \log L}$ --- logarítmico
en el tamaño del flush, lineal en $N$.
\end{complexity}

\section{Finalización vía merge $k$-way}

Después de terminar todos los appends, \code{finalize} produce un
único stream ordenado en disco. Si no hubo flushes el camino es
trivial: ordenar el buffer en memoria, listo. Si no, el motor abre
cursores sobre cada run más un cursor sobre el buffer en memoria
ordenado, y los mergea con un min-heap:

\begin{algorithm}
\caption{\textsc{K-Way-Merge}}\label{alg:kway-merge}
\KwIn{cursores $[c_0, \ldots, c_{k-1}]$, archivo de salida $f$}
construir un min-heap binario $H$ sobre los cursores, clavado en
su elemento cabecera (source\_sid, timestamp, dest\_sid)\;
\While{$H \neq \emptyset$}{
    $c \gets$ \code{H.pop\_min()}\;
    escribir el elemento cabecera de $c$ en $f$\;
    \code{c.advance()}\;
    \If{$c$ no está vacío}{\code{H.push($c$)}}
}
\end{algorithm}

\begin{complexity}
Cada uno de los $N$ elementos se popea y pushea exactamente una
vez del heap, pagando $\Ord{\log k}$ por operación, con total
$\Ord{N \log k}$ en tiempo. El heap tiene a lo sumo $k$ cursores
en cualquier instante, así que el uso de RAM es $\Ord{k}$
independiente de $N$ --- esta es la propiedad que permite al
motor procesar un stream de 28~GB en 16~GB de RAM. El output
mergeado se mapea en memoria (mmap) para lecturas posteriores,
así que \code{get\_edges(s)} sigue siendo $\Ord{d(s)}$.
\end{complexity}

\begin{theorem}[Cota de memoria del spill]
En cualquier momento durante la ingesta el footprint en memoria
de \code{SpillableEdgeStore} es a lo sumo \code{spill\_limit}
$\cdot$ $\code{size\_of::<CompactRelation>}$ más unos pocos KB de
bookkeeping.
\end{theorem}

\begin{proof}
\textsc{Spill-Append} flushea el buffer cuando su longitud
alcanza \code{spill\_limit}; entre flushes el buffer guarda a lo
sumo \code{spill\_limit} elementos, cada uno de 32~B, más el
índice \code{HashMap<StrId, Vec<u32>>} que suma $\Ord{L}$
entradas. Los runs en disco no ocupan RAM hasta mergearse. El
merge propiamente dicho sólo mantiene $k$ cursores en el heap,
acotados por la cantidad de flushes; cada cursor pre-bufferea un
bloque IO chico (unos KB).\qed
\end{proof}

\section{Lecturas}

\begin{algorithm}
\caption{\textsc{Get-Edges}}\label{alg:spill-get-edges}
\KwIn{store $S$, source $s$}
\KwOut{slice de \Compact{} para $s$, vacío si desconocido}
\If{$S.\text{finalized}$ and $S.\text{merged\_mmap}$ es \code{Some}}{
    localizar el offset inicial de $s$ vía el índice en disco
    \Comment*[r]{bsearch \Ord{\log n}}
    \Return slice de largo \code{count} desde el mmap\;
}
\If{$\neg S.\text{finalized}$}{
    \Return posiciones de \code{S.buffer\_index[$s$]}
    materializadas desde \code{S.buffer}
    \Comment*[r]{\Ord{d(s)}}
}
\end{algorithm}

\begin{complexity}
\Ord{\log n + d(s)} en el caso finalizado, \Ord{d(s)} en el no
finalizado (más el trabajo de copiar posiciones en un slice).
\end{complexity}

\begin{inthecode}
\filepath{graph\_hunter\_core/src/spill.rs} --- struct y todos
los algoritmos.\\
\filepath{graph\_hunter\_core/src/spill.rs:327} --- entry point
de \code{finalize}.\\
\filepath{graph\_hunter\_core/src/spill.rs:524} ---
\code{get\_edges}.\\
\filepath{graph\_hunter\_core/benches/dedup\_throughput.rs} ---
bench de ingesta con spill.
\end{inthecode}
